\section{Discussion and limitations}
\label{sec:limits}
\paragraph{What the evidence supports.} In score space and under the stated caps, the aggregation rule changes the minimum targeted investment by more than an order of magnitude: the mean rule rewards one large jump summed over the leaderboard and is expensive in total points, while rules based on order statistics or on counting (median, minimum, plurality) are cheap in total points for near-top targets. The minimum and the power means are exposed to caps (first place becomes unattainable for some targets), and the median's cheapness comes with a large relative spread of cost under evaluation noise (the minimum's is larger still). The common rule of thumb that a median or minimum ``forces breadth'' holds for weak targets (worst case attained: $K=29$ on HELM), but for challengers near the top the median rule's $K$ grows only slowly with the number of tasks and the minimum rule's hardly at all. On the six-task Open LLM Leaderboard the cost of the median rule and the minimum relative to the mean depends on the score scale in paired comparisons (median rule: ranks 11--100; minimum: every group), not for the median rule beyond rank 100.

\paragraph{Where conclusions do not hold or were not tested.}
(1)~\emph{Modelling assumption.} Investment is a score increase within a cap; we make no statement about the real cost of fine-tuning, whether gains transfer, or contamination. empirical caps from fine-tune/base differences (540 pairs) exist only for the Open LLM Leaderboard, and whether a required gain is plausible depends on that envelope assumption, not on the data alone.
(2)~\emph{One-sided.} Rivals do not react; this is not an equilibrium analysis like \citet{chen2026incentives}.
(3)~\emph{Small $n$.} With six tasks (OLLB) the median, minimum and plurality are coarse; rule-versus-$n$ claims come from HELM subsets.
(4)~\emph{Few targets, dependent models.} The primary target set has ten models, so intervals are wide, and models share bases and training data; item-level noise is addressed by resampling and margins, not by modelling dependence.
(5)~\emph{Proxy for pairwise rules.} The Condorcet requirement upper-bounds Copeland, Schulze and Kemeny; their exact costs are not reported.
(6)~\emph{Not completed.} Ordinal rules on the Open LLM Leaderboard ran on a 102-target sample and not at every cap; noise, bootstrap and task-subset robustness ran on HELM only; there is no statistical margin for the rank rules and no exact Copeland (Appendix~\ref{app:notdone}).
